
\documentclass{article}
\usepackage{colortbl}
\usepackage{makecell}
\usepackage{multirow}
\usepackage{supertabular}

\begin{document}

\newcounter{utterance}

\centering \large Interaction Transcript for game `ifeval', experiment `ifeval', episode 144 with SLED{-}BSU\_\_Qwen3.5{-}27B{-}sft{-}ep1\_\_prm{-}ep1{-}best\_of\_n.
\vspace{24pt}

{ \footnotesize  \setcounter{utterance}{1}
\setlength{\tabcolsep}{0pt}
\begin{supertabular}{c@{$\;$}|p{.15\linewidth}@{}p{.15\linewidth}p{.15\linewidth}p{.15\linewidth}p{.15\linewidth}p{.15\linewidth}}
    \# & \multicolumn{2}{c}{Player} && \multicolumn{2}{c}{Game Master} \\
    \hline

    \theutterance \stepcounter{utterance}  
    & & & \multicolumn{4}{p{0.6\linewidth}}{
        \cellcolor[rgb]{0.9,0.9,0.9}{
            \makecell[{{p{\linewidth}}}]{
                \texttt{\tiny{[P1$\langle$GM]}}
                \texttt{The Jimenez family has 20 members. Each member has at least one pet. If there are 32 pets in total, what is the maximum number of members that can possibly have 3 pets? Solve this step by step, and wrap your entire response (reasoning process and the final answer) with double quotation marks.} \\
            }
        }
    }
     \\ \\

    \theutterance \stepcounter{utterance}  
    & \multicolumn{4}{p{0.6\linewidth}}{
        \cellcolor[rgb]{0.561,0.737,0.561}{
            \makecell[{{p{\linewidth}}}]{
                \texttt{\tiny{[P1$\rangle$GM]}}
                \texttt{"To find the maximum number of family members who can have 3 pets, we must minimize the number of pets held by the other members.} \\
\\ 
\texttt{Step 1: Understand the constraints.} \\
\texttt{There are 20 family members.} \\
\texttt{Each member has at least 1 pet.} \\
\texttt{The total number of pets is 32.} \\
\\ 
\texttt{Step 2: Set up the equation.} \\
\texttt{Let \$x\$ be the number of members who have 3 pets.} \\
\texttt{Then, the remaining \$(20 {-} x)\$ members have the minimum possible number of pets, which is 1 pet each.} \\
\\ 
\texttt{Step 3: Formulate the equation for the total number of pets.} \\
\texttt{Total pets = (Number of members with 3 pets × 3) + (Number of members with 1 pet × 1)} \\
\texttt{\$32 = 3x + 1(20 {-} x)\$} \\
\\ 
\texttt{Step 4: Solve for \$x\$.} \\
\texttt{\$32 = 3x + 20 {-} x\$} \\
\texttt{\$32 = 2x + 20\$} \\
\texttt{\$32 {-} 20 = 2x\$} \\
\texttt{\$12 = 2x\$} \\
\texttt{\$x = 6\$} \\
\\ 
\texttt{Step 5: Verify the solution.} \\
\texttt{If 6 members have 3 pets, they account for \$6 \textbackslash{}times 3 = 18\$ pets.} \\
\texttt{The remaining \$20 {-} 6 = 14\$ members have 1 pet each, accounting for \$14 \textbackslash{}times 1 = 14\$ pets.} \\
\texttt{Total pets = \$18 + 14 = 32\$.} \\
\texttt{This matches the given total.} \\
\\ 
\texttt{Therefore, the maximum number of members that can have 3 pets is 6."} \\
            }
        }
    }
    & & \\ \\

    \theutterance \stepcounter{utterance}  
    & & & \multicolumn{2}{p{0.3\linewidth}}{
        \cellcolor[rgb]{0.9,0.9,0.9}{
            \makecell[{{p{\linewidth}}}]{
                \texttt{\tiny{[GM$|$GM]}}
                {'startend:quotation': {}}
            }
        }
    }
    & & \\ \\

    \theutterance \stepcounter{utterance}  
    & & & \multicolumn{2}{p{0.3\linewidth}}{
        \cellcolor[rgb]{0.9,0.9,0.9}{
            \makecell[{{p{\linewidth}}}]{
                \texttt{\tiny{[GM$|$GM]}}
                \texttt{game\_result = WIN} \\
            }
        }
    }
    & & \\ \\

\end{supertabular}
}

\end{document}
